\documentclass[11pt]{article}
\usepackage{amsmath,amssymb,amsthm,mathrsfs}
\usepackage[margin=1in]{geometry}

\newtheorem{theorem}{Theorem}
\newtheorem{proposition}[theorem]{Proposition}
\newtheorem{lemma}[theorem]{Lemma}
\newtheorem{corollary}[theorem]{Corollary}
\theoremstyle{definition}
\newtheorem{remark}[theorem]{Remark}
\newtheorem{conjecture}[theorem]{Conjecture}

\newcommand{\MM}{\overline{\mathcal{M}}}
\newcommand{\Mo}{\mathcal{M}}
\newcommand{\Z}{\mathbb{Z}}
\newcommand{\Q}{\mathbb{Q}}
\newcommand{\sgn}{\mathrm{sgn}}
\newcommand{\Sym}{\mathrm{Sym}}
\newcommand{\DMT}{\mathbf{DMT}}
\newcommand{\DM}{\mathbf{DM}}
\newcommand{\GM}{G_4}

\title{M113 sketch: Extension of Brown--Fonseca \texttt{arXiv:2508.04844}\\
Proposition D.7 to congruence cover $\Gamma_1(4)$ over $\Z[i,\tfrac{1}{2}]$}
\author{ECI / crossed-cosmos --- Phase 7+1}
\date{2026-05-06}

\begin{document}
\maketitle

\begin{abstract}
We sketch the level-$\Gamma_1(4)$ extension of Brown--Fonseca's
Proposition D.7 (mixed Tate property of $\Mo_{1,n}$ for $n\le 3$).
The result is reduced to two specialist sub-claims (B1) and (B2)
concerning, respectively, the $\mu_6$-stratum at the prime $p=3$
and the alternating-$S_n$ projector composed with $\Gamma_1(4)$
level structure. Probability of full proof: $24$--$28\%$.
\end{abstract}

\section{Setup}

\textbf{Notation.}
\begin{itemize}
\item $\MM_{1,n}$: Deligne--Mumford compactification of the moduli stack of stable
       genus-1 curves with $n$ marked points. Smooth proper DM-stack over $\Z$.
\item $\MM_{1,n}(m)$ (Petersen): smooth proper DM-stack over $\Z[1/m]$
       parametrising stable curves with $n$ marked points and $\Gamma(m)$ level structure
       (full level $m$). $\phi(m) = |U(\Z/m\Z)|$ connected components, each defined
       over $\Z[1/m,\zeta_m]$ (symplectic level structure).
\item $\MM_{1,n}^{\Gamma_1(m)} := \MM_{1,n}(B\Z/m\Z)^{\circ,\mathrm{unr}}$:
       the open and closed substack of $\MM_{1,n}(B\Z/m\Z)$ of \emph{connected} admissible
       $\Z/m$-torsors, \emph{unramified} at each marked point. (Pet12 Remark 6.4.)
       For $n=1$ this recovers $\overline{X}_1(m)$.
\item For $m=4$: $\MM_{1,n}^{\Gamma_1(4)}$ is a smooth proper DM-stack over $\Z[1/4]$,
       components over $\Z[1/4, i]\subset \Z[i, 1/2]$.
\end{itemize}

\textbf{Goal.} Show $M(\MM_{1,3}^{\Gamma_1(4)}) \in \DMT(K)$ for $K=\Q(i)$ or
$K=\Z[i,\tfrac{1}{2}]$ (integrally).

\section{The level-$1$ proof (BF25 Appendix D)}

\begin{theorem}[BF25 Theorem 1.6 / Proposition D.7]
$M(\MM_{1,3}) \in \DMT(\Q)$.
\end{theorem}

\noindent\emph{Sketch (BF25 \S D.4).}
Stratify $\MM_{1,n}$ by finite quotients of products $\Mo_{0,r}\times \Mo_{1,s}$
for $3\le r\le n+2$ and $1\le s\le n$:
\[
\MM_{1,n} \;=\; \bigsqcup_\alpha (\Mo_{0,r_\alpha}\times \Mo_{1,s_\alpha})^{H_\alpha}
\]
where $H_\alpha$ is a finite automorphism group of the dual graph stratum $\alpha$.
Each $\Mo_{0,r}$ is a hyperplane-arrangement complement (mixed Tate). Reduce
to mixed Tate of $\Mo_{1,s}^{\mathrm{node}}$ for $s\le 3$. Lemmas D.2 and D.6 give
explicit stratifications of $\Mo_{1,2}^{\mathrm{node}}$ and $\Mo_{1,3}^{\mathrm{node}}$
over $R = \Z[1/6]$. The resulting strata are
weighted projective stacks $\mathcal{P}(4,6)$, $\mathcal{P}(2,3,4)$,
or quotients of affine spaces by $\mu_2, \mu_4, \mu_6$. All are mixed Tate
over $\Q$ (after base change). \qed

\section{Lifting to $\Gamma_1(4)$ via Petersen}

\begin{lemma}[Pet12 Theorem 5.1 + Remark 6.4]
For $m\ge 1$, the stack $\MM_{1,n}^{\Gamma_1(m)}$ is smooth proper Deligne--Mumford
over $\Z[1/m]$. The alternating part of its cohomology coincides with parabolic
$H^1_!(X_1(m), \Sym^{n-1} R^1\pi_*\Q)$.
\end{lemma}

\noindent\emph{Strategy.} Pull back BF25's level-$1$ stratification along the forgetful map
\[
\pi^{\Gamma}\colon \MM_{1,n}^{\Gamma_1(4)} \longrightarrow \MM_{1,n}.
\]
The map is finite étale away from the cusps and elliptic points (j=0, j=1728)
of $\Gamma_1(4)$. Each level-$1$ stratum $S_\alpha$ pulls back to a finite étale cover
$\widetilde{S}_\alpha \to S_\alpha$ with deck transformation group
\[
\GM := (\Z/4\Z)^* \quad\text{or some subquotient acting on $S_\alpha$}.
\]
Since $\Gamma_1(4)/\{\pm 1\}$ acts via the abelian group $\Z/4$ (det map),
the cover is abelian.

\section{Sub-claim (B1): the $\mu_6$ stratum at $p=3$}

The level-$1$ stratum $[Z/\mu_6]$ (in BF25 Lemma D.2) is the punctured cuspidal
cubic $Z = V(g_2, y^2-4x^3-1)$ inside $\mathbb{A}^3 \setminus \{0\}$,
modded by the $\mu_6$-action. Although $Z$ itself is a punctured elliptic curve
(not mixed Tate), BF25 Remark D.3 observes that $M(Z)^{\mu_6}$ \emph{is} mixed Tate.

The cuspidal cubic $y^2 = 4x^3+c$ has $j=0$, hence the elliptic curve
$E_0\colon y^2=4x^3-1$ has CM by $\Z[\zeta_3]$. Over $\Q$:
\[
H^1(E_0;\Q) \otimes \Q(\zeta_3) \;=\; \Q(\zeta_3) \chi \,\oplus\, \Q(\zeta_3) \overline{\chi}
\]
where $\chi$ is a Hecke character. The $\mu_6$-action eigenspaces are
$\zeta_6 = -\zeta_3^{-1}$ on $\chi$ and $\zeta_6^{-1}$ on $\overline{\chi}$.
\textbf{Hence $H^1(E_0)^{\mu_6} = 0$}. So $M(Z)^{\mu_6}$ has only $H^0$ and $H^2$
contributions, both Tate. \qed[level $1$]

\textbf{At level $\Gamma_1(4)$.} Replace $E_0$ with the Galois cover
$\widetilde{E_0} \to E_0$ corresponding to $\Gamma_1(4)$ level structure on $E_0$
(if it exists; $E_0$ has full $4$-torsion over $\Q(\zeta_{12})$ since $\sqrt{-3}\in\Q(\zeta_{12})$).
The $\Gamma_1(4)\cdot\mu_6$ semi-direct product has order $\le 96$.

\begin{conjecture}[Sub-claim (B1)]
$M(\widetilde{Z})^{\Gamma_1(4)\rtimes \mu_6} \in \DMT(\Z[i,\tfrac{1}{2}])$.
\end{conjecture}

\noindent\emph{Plausibility.} The $\Gamma_1(4)$-action on $\widetilde{E_0}$ commutes
with the $\mu_6$ CM action only after base change to $\Q(\zeta_{12})$.
Over $\Q(i)$ with $i\not\in \Q(\zeta_3)$, the semi-direct product is non-trivial.
\textbf{This requires explicit computation.} Most likely outcome (heuristic):
the result holds over $\Z[i,\zeta_3,\tfrac{1}{6}]$, and \emph{descends to $\Z[i, \tfrac{1}{2}]$
modulo a residual issue at $p=3$}.

\section{Sub-claim (B2): alternating-$S_3$ projector at level $\Gamma_1(4)$}

The $S_3$-action on $\MM_{1,3}^{\Gamma_1(4)}$ permutes the $3$ marked points
(but \emph{not} the $\Gamma_1(4)$-distinguished $4$-torsion section).
The cusp form motive is the alternating part:
\[
M^3_{\mathrm{cusp}, \Gamma_1(4)} \;:=\; M(\MM_{1,3}^{\Gamma_1(4)})_\varepsilon[3].
\]
Pet12 Theorem 5.1 + Remark 6.4 imply this realises
$H^1_!(X_1(4), \Sym^2 R^1\pi_*\Q)$, the space of weight-$4$ cusp forms for $\Gamma_1(4)$.

\textbf{But weight-$4$ cusp forms for $\Gamma_1(4)$ are zero!}
(Standard fact: $\dim S_4(\Gamma_1(4)) = 0$.)

So $M^3_{\mathrm{cusp}, \Gamma_1(4)} = 0$ \emph{at the level of Betti realisation}.
By BF25 Theorem 10.13's argument (via Conjecture 1 = conservativity for cusp form
motives), this would imply $M^3_{\mathrm{cusp}, \Gamma_1(4)} = 0$ as a motive.

\begin{proposition}[Implication of (B2)]
If $M^3_{\mathrm{cusp}, \Gamma_1(4)}$ vanishes as a motive (not just in Betti),
then the alternating part of $M(\MM_{1,3}^{\Gamma_1(4)})$ is zero. Combined with
(B1) and the level-$1$ proof, we get $M(\MM_{1,3}^{\Gamma_1(4)}) \in \DMT(\Q(i))$.
\end{proposition}

\textbf{Status of (B2).} BF25 explicitly notes (\S 10.5.2):
\textit{``in every weight, level and nebentypus for which the corresponding space of
cusp forms vanishes, one hopes to prove\ldots\ that the corresponding motive vanishes.''}

This is exactly our situation: weight $4$, level $\Gamma_1(4)$,
$\dim S_4(\Gamma_1(4)) = 0$. So sub-claim (B2) reduces to the
\textbf{conservativity conjecture} (Ayoub Conj 2.12 / Voevodsky standard
conjecture) restricted to cusp form motives at level $\Gamma_1(4)$, weight $4$.

\section{Verdict (B) REDUCED}

\begin{theorem}[Conditional]
Assume:
\begin{enumerate}
\item[(B1)] $M(\widetilde{Z})^{\Gamma_1(4)\rtimes \mu_6} \in \DMT(\Z[i,\tfrac{1}{2}])$
            for the $\mu_6$-cuspidal stratum $Z$ of BF25 Lemma D.2.
\item[(B2)] The cusp form motive $M^3_{\mathrm{cusp}, \Gamma_1(4)}$ vanishes as a
            motive (not just in Betti), \emph{or equivalently} the conservativity
            conjecture for cusp form motives holds in weight $4$, level $\Gamma_1(4)$.
\end{enumerate}
Then $M(\MM_{1,3}^{\Gamma_1(4)}) \in \DMT(\Z[i,\tfrac{1}{2}])$ and the
Gross--Zagier conjecture for $\Gamma_1(4)$, weight $4$, holds for any pair of CM points
defined over $\Q(i)$ (or by extension, $\Z[i,\tfrac{1}{2}]$).
\end{theorem}

\section{ECI implications}

If the conditional theorem above holds:
\begin{enumerate}
\item ECI's M52 ``$6/5$ anchor'' (Sym$^2$($R^1\pi_*$) for $j$-invariant CM tower over $\Q(i)$)
       is realised as a single-valued period of the Brown--Fonseca biextension
       at level $\Gamma_1(4)$.
\item The $\Q(i)$/CM specialisation of M48 ``doublon'' tau structure
       receives a single-valued period interpretation at the M114 cusp.
\item Bayesian update: BF probability M102 22--26\% $\to$ M113 24--28\%.
       Reaching the target 30--40\% requires resolving (B1) at p=3.
\end{enumerate}

\section*{References}

\begin{thebibliography}{99}

\bibitem{BF25} Francis Brown and Tiago J. Fonseca,
``Single-valued periods of meromorphic modular forms and a motivic interpretation
of the Gross--Zagier conjecture'', arXiv:2508.04844 (v2, 18 Aug 2025), 69 pp.

\bibitem{Pet12} Dan Petersen, ``Cusp form motives and admissible $G$-covers'',
\emph{Algebra Number Theory} \textbf{6}(6):1199--1221, 2012; arXiv:1012.1477.

\bibitem{CF06} Caterina Consani and Carel Faber, ``On the cusp form motives in
genus 1 and level 1'', \emph{Adv. Stud. Pure Math.} \textbf{45}:297--314, 2006.

\bibitem{Inc22} Giovanni Inchiostro, ``Moduli of genus one curves with two marked
points as a weighted blow-up'', \emph{Math. Z.} \textbf{302}(3):1905--1925, 2022.

\end{thebibliography}

\end{document}
